
Uncertainty estimation methods for language models—semantic entropy, self-consistency, token variance, verbalized confidence—have been proposed and validated independently, leaving practitioners without systematic guidance on method selection. We conducted controlled experiments on factual question answering benchmarks (NaturalQuestions, TruthfulQA) using Mistral-7B to characterize mechanistic differences between methods. An ablation study with matched computational budgets (K=10 samples) isolates the contribution of semantic clustering: semantic entropy achieves 0.78 AUROC versus 0.69 for an ensemble baseline without clustering, a 0.09 difference representing 13\% relative improvement. Correlation analysis across four methods reveals maximum pairwise correlation of 0.21 (excluding an implementation issue where token variance and self-consistency produced identical outputs with correlation 1.0), indicating that methods capture orthogonal uncertainty dimensions. However, the hypothesis that error types defined by dataset choice (NaturalQuestions for knowledge gaps, TruthfulQA for confident misconceptions) exhibit distinct uncertainty signatures was not supported: observed diversity difference was 0.10 in the opposite direction from prediction (p=0.158, not significant). These findings demonstrate that semantic clustering contributes measurably beyond multiple sampling, that uncertainty methods measure distinct statistical properties, but that dataset-level labels are insufficient for error-type characterization.
